\section{Results}

\subsection{Model calibration}

\subsubsection{Effects of 2,4-D on larvae} 
%
Exposure to 2,4-D resulted in stagnation of growth at 30 mg/L, 
accompanied by a clear delay in metamorphosis. 
We fitted models based on PMoAs $G$, $M$ and $A$, 
and found that only PMoAs $M$ and $A$ can approximate the observations (Figure \ref{fig:mixvalid_metamorphs}). 
PMoA $G$ could not reproduce this pattern, because even a severe effect of growth implies only a slight delay 
of metamorphosis. 
$A$ explained the timing of metamorphosis slightly better, whereas $M$ explained effects on growth better. 
Based on this calibration, only PMoA $G$ was ruled out.

\begin{figure}[!htb]
    \includegraphics[width=\textwidth]{Discoglossus_24D_summaries_PMoA_comparison.png}
    \caption{
        2,4-D model fit to larval data for different physiological modes of action. 
        Error bands indicate 5th to 95th percentiles. 
        Variability in model output reflects individual variability. 
        Non-monotonicity in the fraction of tadpoles at 3 mg/L is the result of a missing value.
        G = decrease in growth efficiency, M = increase in maintenance costs, A = decrease in assimilation efficiency.
    }
    \label{fig:pmoa_comparison_24d}

\end{figure}

\subsubsection{Effects of Flupyradifurone on larvae}

Flupyradifurone caused effects of growth at 100 mg/L, 
but without a clear delay in metamorphosis. 
Based on the calibration in to 2,4-D data, we could rule out PMoAs $M$ and $A$, 
because both imply a delay in metamorphosis when growth is affected. 
Since the timing of metamorphosis was slightly earlier in the affected treatment, 
we also considered PMoA $\kappa$, which causes increased investment in maturation. 
PMoA $G$ clearly explained effects on growth better, while implying a slight delay in metamorphosis 
(Figure \ref{fig:pmoa_comparison_flp}). 
PMoA $\kappa$ over-predicted the effects on growth, causing severy shrinking at high concentrations, 
while also over-predicting the effects on premature metamorphosis. 
Based on this calibration, PMoA $G$ is more plausible.

\begin{figure}[!htb]
    \includegraphics[width=\textwidth]{Discoglossus_Flupyradifurone_summaries_PMoA_comparison.png}
        \caption{
        Flupyradifurone model fit to larval data for different physiological modes of action. 
        Error bands indicate 5th to 95th percentiles. 
        Variability in model output reflects individual variability. 
        G = decrease in growth efficiency, $\kappa$ = decreased resource allocation to growth/increased resource allocation to maturation.
    }
    \label{fig:pmoa_comparison_flp}

\end{figure}

\subsection{Cross-validation: Effects of 2,4-D at GS 46}
2,4-D caused delay of GS 46 and a decrease 
in the mass at GS 46. 
The effect of timing was predicted imprecisely, but accurately. 
The prediction accuracy for the timing of GS 46 was similar for PMoAs $M$ and $A$ (Figure \ref{fig:metamorps_24D}).
Effects on mass at GS 46 were under-predicted by both PMoAs, 
but less so by $M$. 
We  therefore concluded that $M$ is the more plausible PMoA to explain effects of 2,4-D. 
The associated MAPE was 22\% for the timing of GS 46 and 17\% for mass at GS 46. 

\begin{figure}[!htb]
    \includegraphics[width=\textwidth]{CrossValidation_Discoglossus_24D_PMoA_comparison.png}
    \caption{
        Forward predictions of 2,4-D effects on metamorphosis traits based on two alternative PMoAs. 
        Distributions are plotted as kernel density estimates and show the distributions across individuals. 
        Variability in model output reflects individual variability. 
        M = increase in maintenance costs, A = decrease in assimilation efficiency.  
        G46 = developmental stage gosner stage 46.
        MAPE = mean absolute percent error, calculated only for predictions. 
        }    
    \label{fig:metamorps_24D}
\end{figure}

\subsection{Cross-validation: Effects of Flupyradifurone at GS 46}

Effects of Flupyradifurone on the timing of GS 46 were predicted well by PMoA $G$, but not by $\kappa$. 
The latter predicted a premature completion of metamorphosis, which was not observed. 
Effects on mass at GS 46 were predicted accurately by both PMoAs (Figure \ref{fig:metamorps_flupy}). 
Therefore, we concluded 

\begin{figure}[!htb]
    \includegraphics[width=\textwidth]{CrossValidation_Discoglossus_Flupyradifurone_PMoA_comparison.png}
    \caption{
        Forward predictions of Flupyradifurone (Sivanto) effects on metamorphosis traits based on two alternative PMoAs. 
        Distributions are plotted as kernel density estimates and show the distributions across individuals. 
        Variability in model output reflects individual variability. 
        G46 = developmental stage gosner stage 46.
        G = decrease in growth efficiency, $\kappa$ = decreased investment in growth/increased investment in maturation.  
        MAPE = mean absolute percent error = mean relative error, calculated only for predictions. 
        }
    \label{fig:metamorps_flupy}
\end{figure}


\subsection{Model validatiation: Predicting mixture effects}

\subsubsection{Predicting mixture effects on larvae}

Based on the insights from calibration and cross-validation, 
we generated predictions of mixture toxicity by assuming PMoA $M$ for 2,4-D and $G$ for Flupyradifurone. \\
In the mixture treatments, the data indicated a dominance effect of 2,4-D with respect to growth 
and timing of metamorphosis. 
That is, additional exposure to Flupyradifurone did not alter the response to 2,4-D. 
This was predicted accurately by the model (Figure \ref{fig:mixture_growth}, Figure \ref{fig:mixture_fracttadpoles}).
In the highest 2,4-D treatment, the model predicted severe shrinking and a complete inability to reach metamorphosis, 
which is generally in line with the observation that this treamtent caused 100\% mortality.
%
\begin{figure}[!htb]

    \includegraphics[width=\textwidth]{ModelValidation_Discoglossus_UCLM_mixture_growth.png}
    \caption{
        Forward predictions of binary mixture effects on larval growth. 
        Rows correspond to 2,4-D exposure concentrations, 
        columns to Flupyradifurone exposures (both in $mg/L$).
        Variability in model output reflects individual variability. 
        Error bands indicate 5th to 95th percentiles.
    }
    \label{fig:mixture_growth}
    
\end{figure}
%
\begin{figure}[!htb]
    \includegraphics[width=\textwidth]{ModelValidation_Discoglossus_UCLM_mixture_fracttadpoles.png}
    \caption{
        Forward predictions of binary mixture effects on the time to reach metamorphosis (Gosner stage 42).
        \textit{Fraction of tadpoles} is the number of test organisms which have not yet reached metamorphosis, 
        relative to the total number of surviving individuals. \\ 
        Rows correspond to 2,4-D exposure concentrations, 
        columns to Flupyradifurone exposures (both in $mg/L$).
        Variability in model output reflects individual variability. 
        Error bands indicate 5th to 95th percentiles.
    }
    \label{fig:mixture_fracttadpoles}
    
\end{figure}
%
\subsubsection{Predicting mixture effects on metamorphs}
%
Similarly to the result of the cross-validation, the model predicted 
mixture effects on the timing of GS 46 with high accuracy, but low accuracy (Figure \ref{fig:mixture_metamorphs}). 
This coincided with a high mortality at the corresponding treatments. 
The model predictions imply that individuals, if they would not have died from the epxosure, 
may have taken a long time to complete metamorphosis, given that they reach GS 42 at all. 
Remarkably, the predicted variability in the timing of GS 46 was exclusively the result 
of the simulated individual variability, with the degree of variability estimated from control data \citep{hansulDynamicModellingAmphibian}.
This is different from the potential variability in simulation outputs that could result from the propagation of parameter uncertainty. \\
The model under-predicted the effects on mass at GS 46 at the higher mixture treatment. 
When we compare this to the predictions of single-stressor effects, 
it becomes apparent that this mostly reflects a tendency of the model to under-predict effects on mass at GS 46, 
rather than under-predict the mixture effect.

\begin{figure}[!htb]
    \includegraphics[width=\textwidth]{ModelValidation_Discoglossus_UCLM_mixture_metamorphs.png}
    \caption{
        Predicted effects of Flupyradifurone/Sivanto and 2,4-D on metamorphosis traits. 
        Error bands indicate 5th to 95th percentiles. 
        Variability in model output reflects individual variability. 
        G46 = GS 46. FLP = Flupyradifurone.  
    }
    \label{fig:mixture_metamorphs}
\end{figure}